% Configurazione condivisa per gli appunti di Data Mining (Italiano).
\usepackage[T1]{fontenc}
\usepackage[utf8]{inputenc}
\usepackage[italian]{babel}
\usepackage{lmodern}
\usepackage{microtype}

% Mathematics
\usepackage{amsmath,amssymb,amsthm}
\usepackage{mathtools}

% Algorithms: numbered lines and readable syntax
\usepackage[ruled,vlined,linesnumbered]{algorithm2e}
\SetKwInput{KwInput}{Dati}
\SetKwInput{KwOutput}{Risultato}
\SetKwComment{tcp}{ // }{}
\DontPrintSemicolon

% Graphs and diagrams
\usepackage{tikz}
\usetikzlibrary{positioning,arrows.meta,trees,shapes.geometric}

% Tables and layout
\usepackage{booktabs}
\usepackage{tabularx}
\usepackage{array}
\newcolumntype{Y}{>{\raggedright\arraybackslash}X}

% Source code
\usepackage{xcolor}
\usepackage{listings}
\definecolor{codebackground}{HTML}{F5F6F7}
\definecolor{codecomment}{HTML}{64748B}
\definecolor{codekeyword}{HTML}{1D4ED8}
\lstdefinestyle{code}{
  backgroundcolor=\color{codebackground},
  basicstyle=\ttfamily\small,
  commentstyle=\color{codecomment},
  keywordstyle=\color{codekeyword}\bfseries,
  numbers=left,
  numberstyle=\scriptsize\color{codecomment},
  frame=single,
  breaklines=true,
  showstringspaces=false
}

% Visual callout boxes
\usepackage[most]{tcolorbox}
\newtcolorbox{alertblock}[2][]{
  colback=red!5!white,
  colframe=red!75!black,
  fonttitle=\bfseries\small,
  title={#2},
  boxrule=0.8pt,
  arc=3pt,
  #1
}
\newtcolorbox{noteblock}[2][]{
  colback=blue!5!white,
  colframe=blue!75!black,
  fonttitle=\bfseries\small,
  title={#2},
  boxrule=0.8pt,
  arc=3pt,
  #1
}

% Formal environments in Italian
\theoremstyle{plain}
\newtheorem{theorem}{Teorema}[chapter]
\newtheorem{lemma}[theorem]{Lemma}
\newtheorem{property}[theorem]{Propriet\`a}

\theoremstyle{definition}
\newtheorem{definition}[theorem]{Definizione}
\newtheorem{example}[theorem]{Esempio}
\newtheorem{remark}[theorem]{Nota}

% Hyperref
\usepackage[colorlinks=true,hypertexnames=false,linkcolor=blue!50!black,citecolor=blue!50!black,urlcolor=blue!50!black]{hyperref}
\hypersetup{
  pdftitle={Appunti di Data Mining (IT)},
  bookmarksopen=true
}

\newcommand{\BigO}[1]{\mathcal{O}\!\left(#1\right)}
